% =====================================================================
%  CAPÍTULO 1 — INTRODUCCIÓN
% =====================================================================
\chapter{Introducción}
\label{cap:introduccion}

\section{Contexto y justificación}
\label{sec:contexto}

El 5 de enero de 2023 entró en vigor la Directiva (UE) 2022/2464, conocida como
\textit{Corporate Sustainability Reporting Directive} (CSRD), y con ella se abrió la que debía ser una
de las transformaciones más profundas que ha vivido la información corporativa europea en las
últimas décadas. Sin embargo, este cambio no es tan fácilmente perceptible.
El ejercicio 2024 ha sido el primero en el que las grandes empresas europeas de interés público
han estado obligadas a reportar su desempeño en sostenibilidad conforme a un marco
único, los \textit{European Sustainability Reporting Standards} (ESRS), cuyos primeros
informes se han publicado a lo largo de 2025. Por primera vez, miles de compañías han tenido
que aplicar el principio de doble materialidad, cuantificar sus emisiones a lo largo de toda
la cadena de valor y someter parte de esa información a verificación externa. El reporting de
sostenibilidad ha dejado de ser un ejercicio mayoritariamente voluntario y promocional para
convertirse en una obligación regulada, comparable y auditable \parencite{csrd2022}.

Esto se produce en un contexto en el que la información no financiera ha
ganado un peso enorme. Para inversores institucionales, supervisores y consumidores cada vez es más 
importante saber no solo cuánto gana una empresa, sino cómo lo gana: qué impacto ambiental y
social genera, qué riesgos de transición asume y qué compromisos adquiere de cara al futuro.
La sostenibilidad ha pasado a ser fundamental en la comunicación
corporativa, y los informes anuales dedican un espacio creciente a explicarla. Ahora bien,
esa misma centralidad ha hecho aflorar un problema bien conocido en la literatura: la
distancia entre lo que las empresas \textit{dicen} y lo que realmente \textit{hacen}. Lo conocemos como
\textit{greenwashing} (la construcción de una imagen de sostenibilidad más
favorable que el desempeño efectivo) o su versión más sutil, el \textit{cheap talk}
(el compromiso retórico que no se traduce en acciones concretas ni en datos verificables). Este fenómeno
se han convertido en una de las grandes preocupaciones regulatorias y académicas del momento
\parencite{bingler2022}.

Este es el origen de este trabajo. Buena parte de los análisis de
sostenibilidad descansan sobre las calificaciones o \textit{ratings} ESG que elaboran
proveedores externos como MSCI, Sustainalytics o Refinitiv. Estos indicadores son cómodos,
pero presentan dos limitaciones serias: son poco transparentes y muestran
divergencias notables entre proveedores para una misma empresa, haciendo imposible confiar en dichos ratings.
Por ello, el presente trabajo renuncia deliberadamente a cualquier
\textit{rating} de terceros y se construye íntegramente sobre el análisis del \textbf{texto}
de los informes: son las propias palabras de las empresas las que se interrogan.

Para hacerlo a escala, este TFG recurre a técnicas de procesamiento del lenguaje natural
(PLN), una rama de la inteligencia artificial que permite analizar grandes volúmenes de
texto de forma sistemática y reproducible. Este es el eje sobre el que gira el trabajo:
la combinación de un marco regulatorio en plena transición, un problema de relevancia económica 
como el \textit{greenwashing} y un instrumental analítico capaz de leer miles de páginas de informes resulta especialmente
oportuna para conseguir información real y honesta sobre las empresas del EUROSTOXX 600.

\bigskip
La elección del tema responde también a una motivación personal. A lo largo del grado, lo
que más me ha atraído ha sido el punto en el que el análisis de datos se cruza con las
decisiones empresariales, la idea de que detrás de un texto aparentemente cualitativo
se esconden patrones medibles que cuentan una historia.
Cuando se acercó el momento de plantear el TFG me pareció una
oportunidad inmejorable para llevar esa inquietud a un problema económico y financiero real,
de plena actualidad y con implicaciones tangibles para empresas, inversores y reguladores.
Aplicar la inteligencia artificial a un terreno tradicionalmente reservado al análisis manual
era, para mí, mejor manera de acabar mi carrera, poniendo a trabajar juntas las dos cosas que más me interesan.

\bigskip
Este trabajo justifica su novedad por su oportunidad temporal, por su enfoque metodológico,
y por su contribución inédita: ninguna investigación previa,
hasta donde alcanza la revisión realizada, ha analizado el conjunto del STOXX Europe 600 bajo
el marco completo CSRD/ESRS combinando modelado de temas, análisis de sentimiento y detección
de señales de \textit{greenwashing} sobre un corpus propio.

\section{Objetivos}
\label{sec:objetivos}

El \textbf{objetivo general} de este Trabajo Fin de Grado es analizar, mediante técnicas de
procesamiento del lenguaje natural, cómo las grandes empresas europeas del índice STOXX
Europe 600 comunican su estrategia de sostenibilidad en sus informes corporativos bajo el
marco normativo CSRD/ESRS, e identificar en ese discurso patrones comunicativos y señales
textuales de \textit{greenwashing}.

Este objetivo general se concreta en cuatro objetivos específicos, formulados como preguntas
de investigación que estructuran todo el análisis posterior:

\begin{itemize}
  \item \textbf{RQ1 — Temas predominantes.} Identificar qué temas concentran el
  discurso de sostenibilidad de las empresas del STOXX 600, mediante modelado de temas sobre
  el texto de sus informes, y contrastar esos temas con las categorías de los estándares
  ESRS.
  \item \textbf{RQ2 — Diferencias por sector y país.} Determinar si existen diferencias
  estadísticamente significativas en el contenido y el tono del reporting en función del
  sector de actividad y del país de origen de la empresa.
  \item \textbf{RQ3 — Determinantes y señales de \textit{greenwashing}.} Analizar qué
  factores (sector, país, tamaño) se asocian a una mayor especificidad y cobertura
  cuantitativa de la información, y contrastar la hipótesis de que un tono más optimista se
  corresponde con un discurso menos específico, una señal característica del
  \textit{greenwashing}.
  \item \textbf{RQ4 — Evolución NFRD $\rightarrow$ CSRD.} Examinar cómo evoluciona el
  reporting de sostenibilidad entre el régimen de la Directiva sobre información no financiera
  (NFRD, ejercicios 2022 y 2023) y el primer ejercicio obligatorio de la CSRD (2024).
\end{itemize}

Las conclusiones del trabajo (capítulo~\ref{cap:conclusiones}) se construyen de
forma directa sobre cada uno de ellos, de modo que se pueda verificar que todas las
preguntas planteadas al inicio reciben respuesta.

\section{Metodología y fuentes de información}
\label{sec:metodologia-resumen}

El trabajo adopta un enfoque cuantitativo de análisis de contenido textual asistido por
inteligencia artificial. El texto que se analiza pertenece a la subsección de sostenibilidad de cada
informe corporativo. El detalle completo se expone en el capítulo~\ref{cap:metodologia}, aquí se expone 
un resumen sobre las cosas necesarias para entender la generalidad del trabajo.

En cuanto a las \textbf{fuentes de información}, se distingue entre fuentes primarias y
secundarias. La fuente primaria es un corpus documental construido específicamente para esta
investigación: los informes integrados y anuales, en formato PDF, de una muestra de
\textbf{196 empresas} del STOXX Europe 600, correspondientes a los ejercicios 2022, 2023 y
2024. Tras la extracción y depuración del texto, ese corpus quedó conformado por
\textbf{586 documentos} analizables, de los que se aislaron tanto el informe de gestión como
la subsección de sostenibilidad, dando lugar a un corpus de 1.172 unidades de texto
que, una vez segmentadas, suman aproximadamente 245.400 párrafos y cerca de
540.000 frases. Como fuentes secundarias se emplearon los datos financieros de las
empresas —obtenidos de Yahoo Finance \parencite{yfinance2024}—, los textos normativos
europeos \parencite{nfrd2014,csrd2022,esrs2023,taxonomia2020,sfdr2019} y diversos diccionarios
y recursos lingüísticos especializados.

Sobre ese corpus se aplicó una batería de técnicas de PLN: diccionarios
clásicos de análisis textual financiero \parencite{loughran2011}, modelos de lenguaje
recientes especializados en finanzas y clima \parencite{araci2019,webersinke2022} y 
dos métodos complementarios de modelado de temas \parencite{blei2003,grootendorst2022}.
A partir de sus resultados se construyó un índice sintético de \textit{greenwashing} textual
y se realizaron diversos contrastes estadísticos. Todo el proceso se organizó como un flujo
de trabajo reproducible, estructurado en siete fases sucesivas que van de la construcción de
la muestra a la redacción de este documento.

Por útlimo, cabe resaltar que el análisis prescinde por completo de
calificaciones ESG de terceros. Toda medida de calidad, tono o especificidad procede del
propio texto de los informes, lo que preserva la independencia entre la comunicación y el instrumento de medida.

\section{Orden documental}
\label{sec:orden}

El TFG se organiza en seis capítulos. Tras esta \textbf{introducción}, que
delimita el contexto, los objetivos y el enfoque general, el \textbf{capítulo~2} desarrolla el
marco del trabajo. Este se aborda desde dos perspectivas: por un lado, el marco
normativo, que recorre la evolución de la regulación europea de la información de
sostenibilidad desde la NFRD hasta la CSRD, los ESRS, la Taxonomía de la UE y el reglamento
SFDR y, por otro, el marco teórico, que recoge las aportaciones de la literatura sobre
legitimidad organizativa, gestión de impresiones y \textit{greenwashing} y el estado
de la cuestión en la aplicación del PLN al reporting de sostenibilidad.

El \textbf{capítulo~3} detalla la metodología: el diseño de la investigación y las preguntas
que la guían, la construcción de la muestra y del corpus, las técnicas de análisis textual
empleadas, la construcción del índice de \textit{greenwashing} y lso resultados utilizados para contrastar las hipótesis.

El \textbf{capítulo~4} presenta y discute los resultados, organizados en torno a las cuatro
preguntas de investigación: la radiografía del corpus y su cobertura de los estándares ESRS,
los temas predominantes, el tono y la especificidad del discurso, el índice de
\textit{greenwashing}, los factores que lo determinan y, finalmente, la evolución del
reporting entre los dos regímenes normativos.

El \textbf{capítulo~5} recoge, a modo de autocrítica, las principales limitaciones del trabajo
y las propuestas de mejora que de ellas se derivan. El \textbf{capítulo~6} cierra el documento
con las conclusiones y las futuras líneas de investigación. Por último la bibliografía y los anexos
recogen el material complementario que respalda el análisis.

Por último, los resultados se han volcado además en un
\textbf{cuadro de mando interactivo} de elaboración propia, que permite explorar visualmente
el corpus, las empresas individuales, los temas detectados y los principales hallazgos por
pregunta de investigación.
